\documentclass[12pt,a4paper,twocolumn]{article}

\usepackage[utf8]{inputenc}
\usepackage[T1]{fontenc}
\usepackage[polish]{babel}
\usepackage[margin=1.8cm]{geometry}
\usepackage{amsmath}
\usepackage{amssymb}
\usepackage{booktabs}
\usepackage{graphicx}
\usepackage{hyperref}
\usepackage{listings}
\usepackage{xcolor}
\usepackage{caption}
\usepackage{float}
\usepackage{array}
\usepackage{microtype}

\lstset{
  basicstyle=\ttfamily\footnotesize,
  keywordstyle=\color{blue},
  commentstyle=\color{gray},
  frame=single,
  breaklines=true,
  language=Python,
}

\title{\textbf{Autonomiczny samochód wyścigowy\\z użyciem Proximal Policy Optimisation}\\
       \large Projekt -- Głębokie Uczenie ze Wzmocnieniem}
\author{}
\date{}

\begin{document}

\maketitle

%------------------------------------------------------------
\section{Opis problemu}

Celem projektu jest nauczenie agenta autonomicznego pokonywania okrążeń na dwuwymiarowym torze wyścigowym bez żadnej zakodowanej na stałe strategii jazdy.
Środowisko opiera się na symulacji Pygame. Agent otrzymuje wyłącznie informację wizualną -- wycinek obrazu wyśrodkowany na pojeździe -- i uczy się sterowania wyłącznie na podstawie sygnału nagrody.

%------------------------------------------------------------
\section{Środowisko}

\subsection{Symulacja}

Tor jest stałą jednopętlową trasą o rozdzielczości $810 \times 810$ pikseli.
Silnik fizyczny pochodzi z klasy \texttt{AbstractCar}: pojazd ma ciągłą prędkość i kąt obrotu oraz pięć dyskretnych akcji (\textit{przód}, \textit{tył}, \textit{lewo}, \textit{prawo}, \textit{stop}).
Zderzenie z granicą toru wywołuje \texttt{bounce()}, co neguje prędkość.

Aby zmaksymalizować przepustowość treningu, zegar symulacji jest ustawiony na (\texttt{FPS = 10000}).

\subsection{Cykl epizodu}

Epizod zaczyna się, gdy samochód zostaje umieszczony w pozycji startu o współrzędnych $(180, 200)$.
Kończy się -- a pojazd natychmiast wraca na start -- w dwóch przypadkach:

\begin{itemize}
  \item \textbf{Ukończenie okrążenia}: samochód przejeżdża przez linię mety z właściwego kierunku. Zdarzenie to jest odnotowywane jako \emph{sukces}; licznik okrążeń rośnie, a trening jest kontynuowany bez przerwy.
  \item \textbf{Przekroczenie limitu czasu}: pojazd przez $n$ kolejnych kroków nie zalicza żadnego nowego punktu kontrolnego. Zdarzenie jest odnotowywane jako \emph{porażka}. Wartość zmiennej $n$ była zmienna zależnie od etapu treningu, początkowo wynosiła ona 128 potem zwiększoną tą wartość do 1024.
\end{itemize}

Bufor ramek oraz licznik postępu na checkpointach są czyszczone przy każdym resecie gry, tak aby nowy epizod startował z czystego stanu.

%------------------------------------------------------------
\section{Przestrzeń obserwacji}

\subsection{Mapa toru}

Zamiast pikseli RGB środowisko utrzymuje statyczną mapę \texttt{TRACK\_MAP} o kształcie $(H, W)$ (ta sama rozdzielczość co okno gry).
Jest ona obliczana raz przy starcie na podstawie kanału alfa sprite'a krawężnika:

\begin{equation}
  \texttt{TRACK\_MAP}[r,c] =
  \begin{cases}
    0{,}0 & \text{jeśli } \alpha[r,c] > 0 \;\; (\text{ściana}) \\
    0{,}5 & \text{w przeciwnym razie} \;\; (\text{jezdnia})
  \end{cases}
\end{equation}

Przy generowaniu każdej ramki do mapy dokładane są dynamicznie cztery wartości semantyczne:

\begin{center}
\begin{tabular}{@{}lll@{}}
  \toprule
  Wartość & Znaczenie \\
  \midrule
  0,0 & Ściana / poza obrębem toru \\
  0,5 & Jezdnia \\
  0,75 & Kolejne 30 checkpointów \\
  1,0 & Karoseria pojazdu \\
  \bottomrule
\end{tabular}
\end{center}

\subsection{Wycinek egocentryczny}

Z mapy semantycznej wycinany jest obszar $256\times256$ pikseli za pomocą odwrotnej transformacji afinicznej (\texttt{cv2.warpAffine} z flagami \texttt{WARP\_INVERSE\_MAP | INTER\_NEAREST}), tak że:

\begin{itemize}
  \item Kierunek jazdy samochodu zawsze wskazuje \emph{górę} obrazu wyjściowego.
  \item Pojazd jest zakotwiczony blisko \emph{dolnego środka} ramki (wiersz $r_\text{car} = 256 - h/2 - 2$, gdzie $h$ to wysokość sprite'a), co maksymalizuje widoczność toru przed autem.
\end{itemize}

Macierz transformacji ma postać:
\begin{equation}
  M =
  \begin{bmatrix}
    \cos\theta & \sin\theta & t_x \\
    -\sin\theta & \cos\theta & t_y
  \end{bmatrix}
\end{equation}
\begin{equation}
  t_x = c_x - (\cos\theta \cdot 128 + \sin\theta \cdot r_\text{car})
\end{equation}
\begin{equation}
  t_y = c_y - (-\sin\theta \cdot 128 + \cos\theta \cdot r_\text{car})
\end{equation}

gdzie $(c_x, c_y)$ to środek samochodu, a $\theta$ jego kąt obrotu.

Aby nanieść punkt kontrolny ze współrzędnych światowych $(w_x, w_y)$ na obraz wyjściowy, transformacja jest odwracana analitycznie (obrót jest ortogonalny, więc odwrotność to transpozycja):

\begin{equation}
  \begin{pmatrix}\text{col}\\\text{row}\end{pmatrix}
  =
  \begin{bmatrix}
    \cos\theta & -\sin\theta \\
    \sin\theta & \cos\theta
  \end{bmatrix}
  \begin{pmatrix} w_x - t_x \\ w_y - t_y \end{pmatrix}
\end{equation}

Każdy z kolejnych 30 punktów kontrolnych jest rysowany jako wypełnione koło o promieniu 5 pikseli i wartości 0,75.

\subsection{Stos klatek}

Cztery kolejne semantyczne klatki są łączone wzdłuż osi kanałów, dając tensor stanu $s \in \mathbb{R}^{4 \times 256 \times 256}$.
Jeśli dostępnych jest mniej niż cztery klatki (np. na początku epizodu), pierwsza dostępna klatka jest powielana.

%------------------------------------------------------------
\section{Przestrzeń akcji}

Agent wybiera jedną z pięciu dyskretnych akcji w każdym kroku:

\begin{center}
\begin{tabular}{@{}cll@{}}
  \toprule
  Indeks & Nazwa & Efekt \\
  \midrule
  0 & przód  & Przyspieszenie \\
  1 & tył    & Hamowanie / cofanie \\
  2 & lewo   & Skręt w lewo \\
  3 & prawo  & Skręt w prawo \\
  4 & stop   & Stopniowe spowalnianie \\
  \bottomrule
\end{tabular}
\end{center}

%------------------------------------------------------------
\section{Funkcja nagrody}

Skalarna nagroda $r_t$ w każdym kroku składa się z trzech składników.

\paragraph{Nagroda za postęp.}
Iloczyn skalarny wektora kierunku jazdy $\hat{h}$ i jednostkowego wektora $\hat{d}$ wskazującego na kolejny punkt kontrolny, skalowany aktualną prędkością $v$:
\begin{equation}
  r_\text{postęp} = (\hat{h} \cdot \hat{d}) \cdot \frac{v}{8}
\end{equation}
gdzie $\hat{h} = (-\sin\alpha, -\cos\alpha)$ (konwencja Pygame, oś $y$ skierowana w dół), a $\hat{d}$ to znormalizowany wektor od środka samochodu do najbliższego checkpointu.

\paragraph{Kara za bliskość ściany.}
Trzy promienie (przodem, przód-lewo, przód-prawo) mierzą odległość do krawężnika.
Jeśli którakolwiek jest poniżej 15 pikseli, naliczana jest ujemna kara:
\begin{equation}
  r_\text{ściana} = \sum_{d \in \{d_f,\, d_{fl},\, d_{fr}\}}
    \mathbf{1}[d < 15]\cdot\frac{d - 15}{25}
\end{equation}
Uderzenie w ścianę jest odnotowywane do statystyk, jeśli $r_\text{ściana} < 0$.

\paragraph{Kara za czas.}
Stała wartość $-0{,}025$ na krok zniechęca do stania w miejscu:
\begin{equation}
  r_t = r_\text{postęp} + r_\text{ściana} - 0{,}025
\end{equation}

%------------------------------------------------------------
\section{Architektura sieci}

Polityka i funkcja wartości dzielą konwolucyjny kręgosłup (\textit{backbone}) i rozgałęziają się na oddzielne głowice MLP.
Utrzymywane są trzy niezależne optymalizatory Adam (po jednym na podsieci), ale wszystkie wykonują krok na tej samej zagregowanej stracie.

\subsection{Backbone}

\begin{center}
\begin{tabular}{@{}ll@{}}
  \toprule
  Warstwa & Parametry \\
  \midrule
  Conv2d $\to$ GroupNorm $\to$ GELU & $4{\to}32$, $k{=}5$, $s{=}4$ \\
  Conv2d $\to$ GroupNorm $\to$ GELU & $32{\to}64$, $k{=}5$, $s{=}2$ \\
  Conv2d $\to$ GroupNorm $\to$ GELU & $64{\to}96$, $k{=}5$, $s{=}1$ \\
  Conv2d $\to$ GroupNorm $\to$ GELU & $96{\to}96$, $k{=}3$, $s{=}2$ \\
  Flatten & --- \\
  \bottomrule
\end{tabular}
\end{center}

\textbf{GroupNorm(1, C)} traktuje cały map przestrzenny każdego kanału jako jedną grupę, co odpowiada LayerNorm po wymiarach $(C,H,W)$.

\subsection{Głowica aktora}

\[
\begin{aligned}
  &\xrightarrow{\text{LazyLinear}(128)} \text{GELU} \\
  &\xrightarrow{\text{Linear}(64)} \text{GELU} \\
  &\xrightarrow{\text{Linear}(5)} \text{Softmax} \to \pi(a|s)
\end{aligned}
\]

\subsection{Głowica krytyka}

\[
\begin{aligned}
  &\xrightarrow{\text{LazyLinear}(128)} \text{GELU} \\
  &\xrightarrow{\text{Linear}(64)} \text{GELU} \\
  &\xrightarrow{\text{Linear}(1)} \text{V(s)}
\end{aligned}
\]

(\texttt{feature\_dim = 256}, więc rozmiary warstw ukrytych wynoszą kolejno $128$ i $64$.)

\section{Algorytm uczenia: PPO}

\subsection{Zbieranie danych (rollout)}

Agent wchodzi w interakcję ze środowiskiem przez $N = 2048$ kroków przed jakąkolwiek aktualizacją parametrów.
W każdym kroku zapisywana jest krotka $(s_t, a_t, \log\pi_\text{old}(a_t|s_t), r_t, \text{done}_t)$ na CPU.

\subsection{Uogólniona estymacja przewagi (GAE)}

Wartości $V(s_t)$ dla wszystkich stanów rollout są obliczane w jednym wsadowym przebiegu do przodu (bez gradientów).
Wartość bootstrapowa $V(s_N)$ jest szacowana, jeśli ostatni krok nie był terminalny.
Przewagi są obliczane metodą GAE~\cite{schulman2015high}:

\begin{equation}
  \delta_t = r_t + \gamma V(s_{t+1})(1-d_t) - V(s_t)
\end{equation}
\begin{equation}
  \hat{A}_t = \sum_{l=0}^{\infty} (\gamma\lambda)^l \delta_{t+l},
  \quad \gamma = 0{,}99,\;\; \lambda = 0{,}95
\end{equation}

Zwroty to $\hat{R}_t = \hat{A}_t + V(s_t)$.
Przewagi są następnie normalizowane do zerowej średniej i jednostkowej wariancji w obrębie całego rollout.

\subsection{Aktualizacja PPO}

Przez $K = 8$ epok rollout jest tasowany i przetwarzany w mini-batchach po $32$ próbki.
Łączna strata to:
\begin{equation}
  \mathcal{L} = \mathcal{L}_\text{aktor} + \mathcal{L}_\text{krytyk} + \mathcal{L}_\text{entropia}
\end{equation}

\paragraph{Strata aktora (przycięty surogat).}
\begin{equation}
  \mathcal{L}_\text{aktor} = -\mathbb{E}\!\left[
    \min\!\left(
      r_t(\theta)\hat{A}_t,\;
      \operatorname{clip}(r_t(\theta), 1{-}\varepsilon, 1{+}\varepsilon)\hat{A}_t
    \right)
  \right]
\end{equation}
gdzie $r_t(\theta) = \pi_\theta(a_t|s_t)/\pi_\text{old}(a_t|s_t)$ oraz $\varepsilon = 0{,}1$.

\paragraph{Strata krytyka.}
Błąd średniokwadratowy między przewidywanymi wartościami a zwrotami GAE:
\begin{equation}
  \mathcal{L}_\text{krytyk} = \mathbb{E}\!\left[(V_\theta(s_t) - \hat{R}_t)^2\right]
\end{equation}

\paragraph{Bonus entropii.}
\begin{equation}
  \mathcal{L}_\text{entropia} = -\beta_H \,\mathbb{E}[\mathcal{H}(\pi_\theta(\cdot|s_t))],
\end{equation}
Gdzie współczynnik $\beta$ był dostosowywany zależnie od etapu treningu ręcznie (początkowo był wysoki wraz z czasem był zmniejszany). 

\subsection{Optymalizacja}

Trzy niezależne optymalizatory Adam ze wspólnym współczynnikiem uczenia obsługują odpowiednio: backbone, głowicę aktora i głowicę krytyka.
Normy gradientów są przycinane do wartości $3{,}0$ niezależnie dla każdej podsieci. Wartość stałej uczenia była zmienna zależnie od etapu treningu.

\subsection{Zestawienie hiperparametrów}

\begin{center}
\begin{tabular}{@{}ll@{}}
  \toprule
  Hiperparametr & Wartość \\
  \midrule
  Długość rollout $N$ & 2048 \\
  Liczba epok PPO $K$ & 8 \\
  Rozmiar mini-batcha & 32 \\
  $\gamma$ & 0,99 \\
  $\lambda$ & 0,95 \\
  Parametr przycinania $\varepsilon$ & 0,1 \\
  Przycinanie gradientów & 3,0 \\
  Wymiar cech & 256 \\
  Limit bez postępu & 2048 kroków \\
  \bottomrule
\end{tabular}
\end{center}

\end{document}
